\chapter*{Glossaire}
\addcontentsline{toc}{chapter}{Glossaire}
\begin{description}
\item[Affectation] Rattachement d'un point à une station ouverte. À stations
fixées et sans capacité, une station la plus proche minimise son coût d'accès.
\item[Anneau] Cycle qui relie toutes les stations sélectionnées et revient
à son point de départ.
\item[Aspiration] Règle qui autorise un mouvement tabou lorsqu'il améliore
la meilleure valeur globale rencontrée.
\item[Borne inférieure] Valeur qui ne dépasse pas l'optimum d'un problème
de minimisation. Sa comparaison à une solution réalisable renseigne
sur l'incertitude d'optimalité.
\item[Campagne] Ensemble d'exécutions décrites par une configuration commune,
avec plusieurs instances, cardinalités, méthodes et éventuellement graines.
\item[Cas expérimental] Problème fixé par l'instance, $p$ et $\alpha$.
Une méthode peut y être exécutée sur plusieurs graines.
\item[CBC] Solveur de programmation entière du projet COIN-OR,
appelé ici à travers PuLP.
\item[Construction] Méthode rapide produisant une première solution réalisable.
Dans le projet : grille, ajustement de cardinalité, affectation, plus proche
voisin et 2-opt.
\item[Écart-type d'échantillon] Mesure descriptive de la dispersion des coûts
entre graines, calculée avec le dénominateur $m-1$.
\item[Évaluation] Reconstruction et calcul du coût d'un voisin candidat.
Une itération peut en contenir plusieurs.
\item[Flot de connexité] Quantité fictive envoyée depuis la racine pour imposer
que toutes les stations appartiennent à une même composante.
\item[Formulation compacte] Modèle dont le nombre de variables et de contraintes
est polynomial en la taille des données ; cette propriété ne rend pas
sa résolution entière nécessairement rapide.
\item[Graine] Valeur qui initialise le générateur pseudo-aléatoire et permet
de reproduire une suite de tirages dans un contexte logiciel donné.
\item[Heuristique] Méthode qui cherche une bonne solution sans garantir
qu'elle soit optimale.
\item[Itération] Étape de la recherche. Dans la campagne, la locale évalue
un voisin par étape et le tabou jusqu'à 15.
\item[Métaheuristique] Stratégie générale d'exploration qui organise les
mouvements, la mémoire et les règles d'acceptation d'une recherche.
\item[NP-difficulté] Propriété qui traduit ici l'inclusion d'un problème
combinatoire difficile comme cas particulier. Elle n'implique pas que
toutes les instances soient difficiles en pratique.
\item[Optimum prouvé] Solution dont le statut du solveur établit l'optimalité,
sous les tolérances numériques utilisées.
\item[PLNE] Programmation linéaire en nombres entiers. Le modèle du rapport
est mixte, car les flots sont continus et les décisions sont binaires.
\item[PuLP] Bibliothèque Python de construction de modèles linéaires et
entiers et d'interface avec des solveurs.
\item[Recherche locale] Exploration d'un voisinage en retenant des améliorations.
La version du projet échantillonne un échange à chaque itération.
\item[Recherche tabou] Recherche avec mémoire temporaire de certains retours,
qui peut accepter une dégradation et conserve la meilleure solution visitée.
\item[Solution réalisable] Solution qui respecte toutes les contraintes,
quelle que soit sa qualité de coût.
\item[TSPLIB] Bibliothèque d'instances de problèmes de tournées. Le projet
en réutilise les coordonnées sous sa propre convention de distance.
\item[Voisinage] Ensemble des solutions accessibles par une règle de mouvement,
ici le retrait d'une station non imposée et l'ajout d'un autre point.
\item[2-opt] Amélioration du cycle par suppression de deux arêtes et inversion
du segment intermédiaire afin de réduire la longueur.
\end{description}
